\section{AIME I 1989}
        \begin{enumerate}
        	\item (\textit{AIME I 1989}) Compute $\sqrt{(31)(30)(29)(28)+1}$.


 

	\item (\textit{AIME I 1989}) Ten points are marked on a circle. How many distinct convex polygons of three or more sides can be drawn using some (or all) of the ten points as vertices?


 

	\item (\textit{AIME I 1989}) Suppose $n_{}^{}$ is a positive integer and $d_{}^{}$ is a single digit in base 10. Find $n_{}^{}$ if


 $\frac{n}{810}=0.d25d25d25\ldots$

 

	\item (\textit{AIME I 1989}) If $a<b<c<d<e^{}_{}$ are consecutive positive integers such that $b+c+d$ is a perfect square and $a+b+c+d+e^{}_{}$ is a perfect cube, what is the smallest possible value of $c^{}_{}$?


 

	\item (\textit{AIME I 1989}) When a certain biased coin is flipped five times, the probability of getting heads exactly once is not equal to $0$ and is the same as that of getting heads exactly twice. Let $\frac ij^{}_{}$, in lowest terms, be the probability that the coin comes up heads in exactly $3$ out of $5$ flips. Find $i+j^{}_{}$.


 

	\item (\textit{AIME I 1989}) Two skaters, Allie and Billie, are at points $A$ and $B$, respectively, on a flat, frozen lake. The distance between $A$ and $B$ is $100$ meters. Allie leaves $A$ and skates at a speed of $8$ meters per second on a straight line that makes a $60^\circ$ angle with $AB$. At the same time Allie leaves $A$, Billie leaves $B$ at a speed of $7$ meters per second and follows the straight path that produces the earliest possible meeting of the two skaters, given their speeds. How many meters does Allie skate before meeting Billie?


 
                    \begin{figure}[H]
                        \centering
                        \includegraphics[width=0.5\textwidth]{Problem/AIME_I_1989_6_Problem_Figure_1.png}
                    \end{figure}
                    

 

	\item (\textit{AIME I 1989}) If the integer $k$ is added to each of the numbers $36$, $300$, and $596$, one obtains the squares of three consecutive terms of an arithmetic series. Find $k$.


 

	\item (\textit{AIME I 1989}) Assume that $x_1,x_2,\ldots,x_7$ are real numbers such that

 \begin{align*} x_1 + 4x_2 + 9x_3 + 16x_4 + 25x_5 + 36x_6 + 49x_7 &= 1, \\ 4x_1 + 9x_2 + 16x_3 + 25x_4 + 36x_5 + 49x_6 + 64x_7 &= 12, \\ 9x_1 + 16x_2 + 25x_3 + 36x_4 + 49x_5 + 64x_6 + 81x_7 &= 123. \end{align*}

Find the value of $16x_1+25x_2+36x_3+49x_4+64x_5+81x_6+100x_7$.


 

	\item (\textit{AIME I 1989}) One of Euler's conjectures was disproved in the 1960s by three American mathematicians when they showed there was a positive integer such that 
 \[133^5+110^5+84^5+27^5=n^{5}.\] Find the value of $n$.


 

	\item (\textit{AIME I 1989}) Let $a$, $b$, $c$ be the three sides of a triangle, and let $\alpha$, $\beta$, $\gamma$, be the angles opposite them. If $a^2+b^2=1989c^2$, find


 $\frac{\cot \gamma}{\cot \alpha+\cot \beta}$

 

	\item (\textit{AIME I 1989}) A sample of 121 integers is given, each between 1 and 1000 inclusive, with repetitions allowed. The sample has a unique mode (most frequent value). Let $D^{}_{}$ be the difference between the mode and the arithmetic mean of the sample. What is the largest possible value of $\lfloor D^{}_{}\rfloor$? (For real $x^{}_{}$, $\lfloor x^{}_{}\rfloor$ is the greatest integer less than or equal to $x^{}_{}$.)


 

	\item (\textit{AIME I 1989}) Let $ABCD$ be a tetrahedron with $AB=41$, $AC=7$, $AD=18$, $BC=36$, $BD=27$, and $CD=13$, as shown in the figure. Let $d$ be the distance between the midpoints of edges $AB$ and $CD$. Find $d^{2}$.


 
                    \begin{figure}[H]
                        \centering
                        \includegraphics[width=0.5\textwidth]{Problem/AIME_I_1989_12_Problem_Figure_1.png}
                    \end{figure}
                    

	\item (\textit{AIME I 1989}) Let $S^{}_{}$ be a subset of $\{1,2,3^{}_{},\ldots,1989\}$ such that no two members of $S^{}_{}$ differ by $4$ or $7$. What is the largest number of elements $S^{}_{}$ can have?


 

	\item (\textit{AIME I 1989}) Given a positive integer $n^{}_{}$, it can be shown that every complex number of the form $r+si^{}_{}$, where $r^{}_{}$ and $s^{}_{}$ are integers, can be uniquely expressed in the base $-n+i^{}_{}$ using the integers $0,1,2^{}_{},\ldots,n^2$ as digits. That is, the equation


 $r+si=a_m(-n+i)^m+a_{m-1}(-n+i)^{m-1}+\cdots +a_1(-n+i)+a_0$

 is true for a unique choice of non-negative integer $m^{}_{}$ and digits $a_0,a_1^{},\ldots,a_m$ chosen from the set $\{0^{}_{},1,2,\ldots,n^2\}$, with $a_m\ne 0^{}{}$. We write


 $r+si=(a_ma_{m-1}\ldots a_1a_0)_{-n+i}$

 to denote the base $-n+i^{}_{}$ expansion of $r+si^{}_{}$. There are only finitely many integers $k+0i^{}_{}$ that have four-digit expansions


 $k=(a_3a_2a_1a_0)_{-3+i^{}_{}}~~~~a_3\ne 0.$

 Find the sum of all such $k^{}_{}$.


 

	\item (\textit{AIME I 1989}) Point $P$ is inside $\triangle ABC$. Line segments $APD$, $BPE$, and $CPF$ are drawn with $D$ on $BC$, $E$ on $AC$, and $F$ on $AB$ (see the figure below). Given that $AP=6$, $BP=9$, $PD=6$, $PE=3$, and $CF=20$, find the area of $\triangle ABC$.


 
                    \begin{figure}[H]
                        \centering
                        \includegraphics[width=0.5\textwidth]{Problem/AIME_I_1989_15_Problem_Figure_1.png}
                    \end{figure}
                    

\end{enumerate}